%% kapa: exam=1 round=19 year=2008 period=1 form=A subject=accounting subject_ko=회계학 seq=23 number=103
%% source: Q-Net 감정평가사 기출문제 (cst00309) artlSeq=1001962 file=extracted/1008/1001962/2008년 제19회 감정평가사 1차시험 기출문제/1교시.hwp
%% conversion: kordoc 4.24.51 -> tools/gen1.py (auto); review=proofread-sonnet; flags=none
\begin{question}{103}
(주)계룡에는 A와 B 두 개의 사업부가 있는데 다음은 성과평가와 관련된 자료이다.

{\small\centering
\begin{tabular}{lrr}
\hline
\multicolumn{1}{c}{\uline{구분}} & \multicolumn{1}{c}{\uline{A사업부}} & \multicolumn{1}{c}{\uline{B사업부}} \\
투자액 & 2,000억원 & 4,000억원 \\
영업이익 & 400억원 & 720억원 \\
\hline
\end{tabular}
\par}

(주)계룡의 자본비용은 10\%이며, 각 사업부에 대한 요구수익률도 10\%이다. (주)계룡이 투자수익률과 잔여이익으로 사업부를 평가하는 경우 어떤 결과가 나타나는지 가장 옳은 것은? \par
\begin{choices}
  \item 두 평가방법 모두 A사업부가 더 우수하다.
  \item 두 평가방법 모두 B사업부가 더 우수하다.
  \item 두 평가방법 모두 A와 B사업부간의 성과에는 차이가 없다.
  \item 투자수익률로 평가하는 경우에는 B사업부, 잔여이익으로 평가하는 경우에는 A사업부가 더 우수하다.
  \item 투자수익률로 평가하는 경우에는 A사업부, 잔여이익으로 평가하는 경우에는 B사업부가 더 우수하다.
\end{choices}
\end{question}
\answer{5}
